\section{Impact}\label{sec:impact}

\HECriterion{Impact --- aspects to be taken into account}{%
\begin{itemize}
  \item Credibility of the pathways to achieve the expected outcomes and
  impacts specified in the work programme.
  \item Suitability and quality of the measures to maximise expected outcomes
  and impacts, as set out in the dissemination and exploitation plan, including
  communication activities.
\end{itemize}}

\begin{heguidance}
Show how the project will contribute to the topic's expected outcomes over the
medium term and to the wider expected impacts in the relevant destination over
the longer term. Explain the measures that will maximise those contributions.
\end{heguidance}

\subsection{Project's pathways towards impact\HEBudget{e.g. 3 pages}}
\label{sec:pathways}
\begin{heguidance}
Provide a narrative explaining how the project results will make a difference
beyond the immediate scope and duration of the project. Tailor it to this
project and include the following elements.
\begin{enumerate}[label=(\alph*)]
  \item Describe the unique contribution of the results to the outcomes
  specified in the topic and to the longer-term impacts specified in the
  relevant work programme destination. Give quantified estimates where
  possible and meaningful. Be specific about this project's effects. Identify
  the target groups that benefit, breaking general groups into the relevant
  interest groups or segments of society.
  \item Describe requirements and potential barriers beyond the scope and
  duration of the project that may determine whether the outcomes and impacts
  are achieved. These may concern other R\&I, the regulatory environment,
  targeted markets or user behaviour. Explain how they could evolve and what
  mitigating measures, within or beyond the project, would address barriers or
  incorrect assumptions. Critical project-management risks belong in section
  3.1 and table 3.1e.
\end{enumerate}
\end{heguidance}

\HEParagraph{Contribution to expected outcomes and impacts}
\HEPlaceholder{Map the project results to the topic's expected outcomes and the
destination's expected impacts, identifying beneficiaries and justified
quantitative estimates.}

\HEParagraph{Requirements, barriers and mitigating measures}
\HEPlaceholder{Describe the external assumptions, dependencies, barriers and
measures needed for the proposed pathways to succeed.}

\subsection{Measures to maximise impact --- Dissemination, exploitation and
communication\HEBudget{e.g. 3 pages, excluding section 2.3}}
\label{sec:measures}
\begin{heguidance}
Provide a first version of the plan for dissemination and exploitation including
communication activities. Describe the planned measures and the target groups
addressed, such as the scientific community, end users, financial actors or the
public. This plan is an admissibility condition unless the topic explicitly
states otherwise. A more detailed plan is required as a project deliverable
within 6 months after signature; update it periodically as the project develops.

Communication should promote the project throughout its lifetime, explain the
activities and their use and benefits for citizens, and reach society. State
clear objectives, main messages, tools and channels for each target group.
Activities should start at the outset and continue throughout the project.

Make measures proportionate to the project's scale and specify concrete actions
during and after the project, including standardisation where relevant. Explain
why each measure suits its target group, consider follow-up after completion,
and describe a plausible path to commercialisation where relevant. If
exploitation is expected primarily in non-associated third countries, explain
how it remains in the Union's interest. Describe feedback to policy measures
that could inform existing or new policies and programmes.
\end{heguidance}

\HEParagraph{Dissemination, exploitation and communication plan}
\HEPlaceholder{For each result and target group, specify the intended action,
responsible participant, timing, channel, indicator and follow-up after the
project. Distinguish dissemination, exploitation and communication.}

\HEParagraph{Intellectual property and access to results}
\begin{heguidance}
Outline the intellectual property management strategy, foreseen protection
measures (such as patents, design rights, copyright or trade secrets), and how
these support exploitation. A funded project needs an appropriate consortium
agreement covering ownership of and access to key knowledge, including IPR and
research data. Where relevant, it should enable collective and individual
exploitation of market opportunities. The final periodic report must identify
the owners of results in a results ownership list.
\end{heguidance}
\HEPlaceholder{Explain ownership, protection, access rights, licensing,
exploitation responsibilities and the treatment of jointly developed results.}

\ifHEImpactSummary
\begin{HELandscape}
\subsection{Summary\HEBudget{optional; remove this section if not used}}
\label{sec:impact-summary}
\begin{heguidance}
Summarise the key elements of the impact section in the canvas below. The
summary is optional in template V5.1 and counts towards the overall page limit.
\end{heguidance}
\noindent\textbf{KEY ELEMENT OF THE IMPACT SECTION}\par\medskip
\begin{tcbraster}[raster columns=3,raster equal height=rows,
  raster column skip=4mm,raster row skip=4mm]
  \HECanvasCell{SPECIFIC NEEDS}
    {What are the specific needs that triggered this project?}
    {\HEPlaceholder{Identify the needs addressed by the project.}}
  \HECanvasCell{EXPECTED RESULTS}
    {What do you expect to generate by the end of the project?}
    {\HEPlaceholder{List the main results the project will produce.}}
  \HECanvasCell{D \& E \& C MEASURES}
    {What dissemination, exploitation and communication measures will you
    apply to the results?}
    {\HEPlaceholder{State the measures for uptake, use and communication.}}
  \HECanvasCell{TARGET GROUPS}
    {Who will use or further take up the project results? Who will benefit
    from the results?}
    {\HEPlaceholder{Identify users, adopters and beneficiaries.}}
  \HECanvasCell{OUTCOMES}
    {What change do you expect after successful dissemination and
    exploitation of the results to the target groups?}
    {\HEPlaceholder{Describe the expected changes in use, behaviour or
    practice.}}
  \HECanvasCell{IMPACTS}
    {What wider scientific, economic and societal effects will contribute to
    the expected impacts in the relevant work programme destination?}
    {\HEPlaceholder{State the expected longer-term effects.}}
\end{tcbraster}
\end{HELandscape}
\fi

